\section{Compatibility of the complete common coordinate of two}
\label{sec:eighth-two-compatibility}

Both the eighth-root and the cubic coordinates of $2$ in the
radical--CM intersection are compatible with the corrected CM
frames.  These assertions hold on the multiplier-one gentle group.
We first identify the harmonic torsion level by an exact conductor
computation and prove its congruence directly from the cycle
equations.  An abelianized calculation in the modular group then
gives the cubic coordinate.

\subsection{The relevant Gaussian ray fields}

Put $k=\mathbb Q^{\mathrm{ab}}$, $K_4=\mathbb Q(i)$,
$\mathfrak p=(1+i)$, and $E_4:Y^2=X^3-X$.

\begin{lemma}\label{lem:eighth-two-conductor}
For $n=2,4,8$, the field
$L_n=\mathbb Q(i,2^{1/n})$ is cyclic of degree $n$ over $K_4$.
Its conductor is respectively
$\mathfrak p^4,\mathfrak p^6,\mathfrak p^8$.
In particular
\[
 L_4\subset K_4(E_4[8]),\qquad
 L_8\subset K_4(E_4[16]).
\]
\end{lemma}

\begin{proof}
Write $n=2^r$, $\alpha=2^{1/n}$, and $F_n=\mathbb Q(\alpha)$.
For $n=2,4$, cyclicity follows because $K_4$ contains the required
roots of unity.  For $n=8$ the field $L_8$ contains
$\zeta_8=(1+i)/\alpha^4$, and hence is the splitting field of
$X^8-2$.  Automorphisms fixing $i$ are indexed by
$a\in\mathbb Z/8$, with $\alpha\mapsto\zeta_8^a\alpha$.
Their composition is
$a*b=a+b+4ab$ modulo eight.  The element indexed by $1$ has
order eight, so the group is cyclic.  The degrees are $n$:
$X^n-2$ is Eisenstein over $\mathbb Q$ and $F_n$ is real,
so $[F_n(i):K_4]=n$.

We compute the discriminants without assuming a ray conductor.
The polynomial $X^n-2$ is Eisenstein at two, and $\alpha$ is a
uniformizer of its totally ramified local extension.  Therefore
the local integer ring is $\mathbb Z_2[\alpha]$.
The polynomial discriminant is supported only at two, so
$\mathcal O_{F_n}=\mathbb Z[\alpha]$ globally, and
\begin{equation}
 |\operatorname{disc}(F_n/\mathbb Q)|
       =n^n2^{n-1}=2^{rn+n-1}.
 \label{eq:eighth-two-pure-discriminant}
\end{equation}

In the completion of $F_n$ at its unique prime over two, use
the valuation with $v(\alpha)=1$.  Define
\[
 a_2=1+\alpha,\qquad
 a_4=1+\alpha^2+\alpha^3,\qquad
 a_8=1+\alpha^4+\alpha^6+\alpha^7,
\]
where in each expression $\alpha^n=2$ for its indicated $n$.
Direct expansion gives
\[
 a_n^2+1\equiv\alpha^{2n-1}\pmod{\alpha^{2n}},
 \qquad v(a_n^2+1)=2n-1.
\]
The element
$b_n=(i+a_n)/\alpha^{n-1}$ satisfies
\begin{equation}
 b_n^2-a_n\alpha\,b_n+
              \frac{a_n^2+1}{\alpha^{2n-2}}=0.
 \label{eq:eighth-two-eisenstein}
\end{equation}
Both nonleading coefficients have valuation one, so this is an
Eisenstein quadratic polynomial.  It generates the relative local
integer ring.  Its discriminant has valuation two: the square
of its linear coefficient has valuation two, whereas four times
the constant coefficient has valuation $2n+1$.
Thus the relative discriminant of $L_n/F_n$ has exponent two
at the unique prime over two and has no odd prime factor.
The discriminant tower formula and
\eqref{eq:eighth-two-pure-discriminant} give
\[
 |\operatorname{disc}(L_n/\mathbb Q)|
       =2^{\,2(rn+n-1)+2}=2^{\,2n(r+1)}.
\]
Since $|\operatorname{disc}(K_4/\mathbb Q)|=4$ and
$N_{K_4/\mathbb Q}(\mathfrak p)=2$, a second application of
the tower formula yields
\begin{equation}
 \operatorname{disc}(L_n/K_4)=\mathfrak p^{\,2nr}.
 \label{eq:eighth-two-gaussian-discriminant}
\end{equation}

Let $c_2,c_4,c_8$ be the conductor exponents of faithful characters
of the three cyclic extensions.  Characters of the same order
have the same kernel and hence the same conductor.  The cyclic
tower $L_2\subset L_4\subset L_8$ and the conductor--discriminant
formula \cite[V, Theorem~3.27 and Remark~3.28, p.~165]{MilneCFT}
give
\[
 c_2=4,\qquad c_2+2c_4=16,\qquad
 c_2+2c_4+4c_8=48.
\]
Consequently $(c_2,c_4,c_8)=(4,6,8)$.
The conductors of the extensions themselves are the same as
those of their faithful characters.

The torsion reciprocity theorem for $E_4$ identifies its
torsion points modulo geometric automorphisms with the
corresponding Gaussian ray class fields
\cite[Theorem~12.27, printed p.~124]{MilneMF110}.
Thus the ray class field modulo $(8)=\mathfrak p^6$ is contained
in $K_4(E_4[8])$, and that modulo $(16)=\mathfrak p^8$ is
contained in $K_4(E_4[16])$.  The conductors just computed
give the two asserted inclusions.
\end{proof}

\begin{proposition}\label{prop:eighth-two-division-fields}
Over $k=\mathbb Q^{\mathrm{ab}}$ one has
\[
 k(E_4[8])=k(2^{1/4}),\qquad
 k(E_4[16])=k(2^{1/8}).
\]
\end{proposition}

\begin{proof}
The arithmetic CM image on $G_k$ is
$N_4(4)=\{w\in\mathbb Z_2[i]^\times:
w\bar w=1,\ w\equiv1\pmod4\}$ at the prime two.
This is the image calculated in the adelic frame comparison.
Write $w=a+bi$, where $a=1+4s$ and $b=4t$.
The norm equation gives
$s+2(s^2+t^2)=0$.  In particular $s$ is even and
$s\equiv2t\pmod4$, so
\begin{equation}
 w\equiv1+8t+4ti\pmod{16}.
 \label{eq:eighth-two-norm-reduction}
\end{equation}
Every class $t\bmod4$ occurs: the norm-one unit
$(1+2ti)/(1-2ti)$ has the required imaginary coefficient modulo
sixteen.  Thus the image modulo sixteen is cyclic of order four,
generated by $9+4i$, and its image modulo eight has order two.
It follows that
$[k(E_4[16]):k]=4$ and $[k(E_4[8]):k]=2$.
The two inclusions from Lemma~\ref{lem:eighth-two-conductor}
become inclusions of extensions of these same degrees, since
the radical Kummer image over $k$ is $2\widehat{\mathbb Z}$.
They are therefore equalities.
\end{proof}

\subsection{The cycle equations fix the harmonic frame modulo sixteen}

Use the signed matrices
\[
 U=\begin{pmatrix}1&1\\0&1\end{pmatrix},\quad
 L=\begin{pmatrix}1&0\\-1&1\end{pmatrix},\quad S=ULU,
 \qquad q(x)=U^2,\quad q(y)=L^2.
\]
Write $\gamma_3F=\overline{[F,[F,F]]}$ for the third lower
central subgroup of $F=\widehat F(x,y)$.

\begin{lemma}\label{lem:eighth-two-class-two-image}
Modulo sixteen the commutator subgroup of $q(F)$ is the central
cyclic group of order four generated by
\[
 C=[U^2,L^2]
   =\begin{pmatrix}13&8\\8&5\end{pmatrix}\pmod{16}.
\]
In particular $q_{16}(\gamma_3F)=1$.
\end{lemma}

\begin{proof}
Direct multiplication gives $C^2=9I$, $C^4=I$, and shows that
$C$ commutes with both $U^2$ and $L^2$ modulo sixteen.
In a group generated by two elements, the commutator subgroup
is the normal closure of their commutator.  Here that normal
closure is precisely the central group $\langle C\rangle$.
\end{proof}

\begin{lemma}\label{lem:eighth-two-normalized-gentle}
If $h=(1,f)$ is gentle and $\rho_2(h)=0$, then
\[
 \beta_h\in\gamma_3F,\qquad
 q_{16}(\beta_h)=D_h\bmod16=W_h\bmod16=I.
\]
\end{lemma}

\begin{proof}
In $F/\gamma_3F$, write $f=[x,y]^a$.
At multiplier one, the three-cycle equation makes $3a=0$,
because its three commutators have the same central image.
The commutator coordinate is in the torsion-free group
$\widehat{\mathbb Z}$, so $a=0$ and $f\in\gamma_3F$.
The equation $\rho_2(h)=0$ gives $\beta_h\in F'$.
Write its image modulo $\gamma_3F$ as $[x,y]^d$.
The interchange $\theta(x)=y$, $\theta(y)=x$ inverts this
central commutator, and the frame identity
$f=\theta(\beta_h)\beta_h^{-1}$ gives $2d=0$.
Therefore $d=0$, proving $\beta_h\in\gamma_3F$.

Lemma~\ref{lem:eighth-two-class-two-image} gives
$q_{16}(\beta_h)=I$.  The corrected frame
$W_h=q(\beta_h)^{-1}D_h$ centralizes $S$ as a signed matrix.
Consequently the upper-unipotent matrix $D_h\bmod16$ centralizes
$S$.  Comparing the two diagonal entries in $D_hS=SD_h$
forces its shear to be zero modulo sixteen.  Thus both $D_h$
and $W_h$ are the identity modulo sixteen.
\end{proof}

\begin{theorem}\label{thm:eighth-two-compatibility}
Let $M_0=\ker\lambda$ in the gentle group.  Define the continuous
homomorphism
$\nu_2:N_4(4)\to\mathbb Z/4$ by its arithmetic normalization:
\[
 \nu_2(\rho_4(\sigma))
       =\frac{\kappa_2(\sigma)}2\pmod4
       \qquad(\sigma\in G_{\mathbb Q^{\mathrm{ab}}}).
\]
This is well defined and has kernel the subgroup congruent to
one modulo sixteen.  For every $h\in M_0$,
\begin{equation}
 \nu_2(W_h)=\frac{\rho_2(h)}2\pmod4.
 \label{eq:eighth-two-frame-kummer}
\end{equation}
In particular the exceptional eighth-root coordinate of $2$
in the radical--CM compatibility defect is identically zero
on full GT.
\end{theorem}

\begin{proof}
Proposition~\ref{prop:eighth-two-division-fields} identifies
the quotient of the arithmetic harmonic image modulo sixteen
with $\operatorname{Gal}(k(2^{1/8})/k)$.
The Kummer action on the eighth root has even exponent on $G_k$,
so division by two identifies this group with $\mathbb Z/4$.
This proves the stated definition and kernel of $\nu_2$.

The normalization of the gentle group gives
$M_0=N_0G_k$, where $N_0=\ker(\lambda,\rho_2)$.
Explicitly, $\rho_2(h)$ is even, and the Kummer map
$G_k\to2\widehat{\mathbb Z}$ is onto; choose $\sigma\in G_k$
with $\kappa_2(\sigma)=\rho_2(h)$ and put $n=h\sigma^{-1}$.
Then $n\in N_0$.
By Lemma~\ref{lem:eighth-two-normalized-gentle},
$W_n\equiv I\pmod{16}$.
The corrected-frame homomorphism and its exact arithmetic
comparison give
\[
 \nu_2(W_h)=\nu_2(W_\sigma)
           =\nu_2(\rho_4(\sigma))
           =\kappa_2(\sigma)/2
           =\rho_2(h)/2\pmod4.
\]
This proves \eqref{eq:eighth-two-frame-kummer}.
It is the required compatibility on the multiplier-one subgroup.
Matching an arbitrary multiplier by a Galois element extends the
vanishing to the corresponding crossed-cocycle component on $G$.
\end{proof}

\subsection{The cubic coordinate of two}

Recall the modular presentation
$\Gamma=\langle s,r\mid s^2=r^3=1\rangle$ and its pure subgroup
$F=\ker(\Gamma\to S_3)$, with
$x=u^2$, $y=v^2$, $s=uvu$, $r=uv$, and $sr=u^{-1}$.
For $h\in M_0$ the positive peripheral automorphism satisfies
\[
 \Phi_h(s)=\beta_hs\beta_h^{-1},\qquad
 \Phi_h(r)=A_hrA_h^{-1},\qquad
 \Phi_h(sr)=sr,
 \qquad A_h,\beta_h\in F .
\]
Its corrected cubic frame is
\[
 T_h=q(A_h)^{-1}D_h,\qquad
 R=-UL=\begin{pmatrix}0&-1\\1&-1\end{pmatrix}.
\]
Thus $T_h$ centralizes $R$ as a signed matrix.

\begin{lemma}\label{lem:cubic-two-pure-frame}
For every $h\in M_0$, one has $A_h\in F'$.
\end{lemma}

\begin{proof}
Write the abelian group $F^{\mathrm{ab}}$ additively.
Expanding $\Phi_h(s)\Phi_h(r)=sr$ in the extension
$\Gamma/F'$ gives
\[
 (1-s)[\beta_h]+s(1-r)[A_h]=0.
\]
The harmonic normalization is
$[\beta_h]=-\rho_2(h)([x]+[y])$.
Conjugation by $s$ interchanges $[x]$ and $[y]$, so the first
term vanishes and $(1-r)[A_h]=0$.
Conjugation by $r$ sends
$x\mapsto y$ and $y\mapsto(xy)^{-1}$; its matrix on
$F^{\mathrm{ab}}=\widehat{\mathbb Z}^2$ is therefore
\[
 \begin{pmatrix}0&-1\\1&-1\end{pmatrix}.
\]
The determinant of $1-r$ is $3$.  Since multiplication by $3$
is injective on $\widehat{\mathbb Z}^2$, the adjugate identity
shows that $1-r$ is injective.  Hence $[A_h]=0$.
\end{proof}

The centralizer of $R$ in $\mathrm{SL}_2(\mathbb F_3)$ is
$\{\pm R^j:j\in\mathbb Z/3\}$.  Define
\[
 j_3:N_3(4)\longrightarrow\mathbb Z/3,\qquad
 j_3(t)=j\quad\hbox{if }t\bmod3=(-I)^eR^j .
\]
This is the quotient that kills the central sign.

\begin{proposition}\label{prop:cubic-two-compatibility}
For every $h\in M_0$,
\[
 j_3(T_h)=\rho_2(h)\pmod3.
\]
On the arithmetic subgroup, this quotient is the Kummer action
on $k(2^{1/3})/k$, with the indicated choice of exponent.
\end{proposition}

\begin{proof}
The abelianization of $\mathrm{SL}_2(\mathbb F_3)$ is cyclic of
order three.  Its quotient map $\ell$ can be normalized by
$\ell(U)=\ell(L)=1$; then
$\ell(S)=0$ and $\ell(R)=2$.
For completeness, its derived subgroup is the quaternion group
generated by $S$ and $USU^{-1}$.
These generators have square $-I$, anticommute, and generate
a normal subgroup of order eight.  Conjugation by $U$ permutes
its three cyclic subgroups of order four.  Commutators with $U$
generate this quaternion group, proving the abelianization
assertion.

Write $D_h=\left(\begin{smallmatrix}1&b_h\\0&1\end{smallmatrix}\right)$.
Since $q(x)=U^2$ and $q(y)=L^2$, the harmonic normalization gives
\[
 \ell(q(\beta_h))=-4\rho_2(h)=-\rho_2(h)\pmod3.
\]
The centralizer of $S$ in $\mathrm{SL}_2(\mathbb F_3)$ is
$\langle S\rangle$, so $\ell(W_h)=0$.  The identity
$W_h=q(\beta_h)^{-1}D_h$ now gives
$b_h=-\rho_2(h)$ modulo three.
By Lemma~\ref{lem:cubic-two-pure-frame},
$\ell(q(A_h))=0$, and consequently
\[
 2j_3(T_h)=\ell(T_h)=\ell(D_h)
       =-\rho_2(h)\pmod3 .
\]
Multiplying by $2$ yields the stated formula.

The arithmetic identification can also be seen directly from
$E_3:Y^2=X^3+1$.  The duplication formula
$X(2P)=9X(P)^4/(4Y(P)^2)-2X(P)$ shows that
its three-division polynomial is
$3X(X^3+4)$.  The nonzero three-torsion coordinates satisfy
either $X=0$, $Y=\pm1$, or
$X^3=-4$, $Y^2=-3$.
Since $k$ contains the roots of unity and $\sqrt{-3}$, this gives
\[
 k(E_3[3])=k(2^{1/3}).
\]
The signed arithmetic comparison is $T_\sigma=\kappa(\sigma)
\rho_3(\sigma)$, where $\kappa(\sigma)\in\{\pm1\}$.
Its reduction modulo three can therefore have an additional
central sign.  Passing to $j_3$ kills that sign, and gives the
cubic field just identified; no equality with the full corrected
modulo-three field is asserted.
On $\sigma\in G_k$, the formula already proved reads
$j_3(T_\sigma)=\kappa_2(\sigma)\pmod3$.
This fixes its normalization as the usual Kummer exponent.
\end{proof}

\begin{corollary}\label{cor:complete-two-compatibility}
There is a continuous homomorphism
\[
 \nu_{2,\mathrm{CM}}:N_3(4)\times N_4(4)
          \longrightarrow\mathbb Z/12
\]
specified by
\[
 \nu_{2,\mathrm{CM}}(t,w)\equiv\nu_2(w)\pmod4,\qquad
 \nu_{2,\mathrm{CM}}(t,w)\equiv2j_3(t)\pmod3 .
\]
For every $h\in M_0$,
\[
 \nu_{2,\mathrm{CM}}(T_h,W_h)=\frac{\rho_2(h)}2\pmod{12}.
\]
Arithmetically this is the common quotient
$\operatorname{Gal}(k(2^{1/24})/k)$, with Kummer exponent
$\kappa_2/2$.
In particular the entire order-twelve coordinate belonging to
$2$ in the radical--CM compatibility defect vanishes on full GT.
\end{corollary}

\begin{proof}
The Chinese remainder theorem defines the homomorphism.
The two preceding results give its asserted value because
$2$ is its own inverse modulo three.
The fields $k(2^{1/8})$ and $k(2^{1/3})$ have degrees four and
three over $k$, and together generate $k(2^{1/24})$.
Their arithmetic Kummer exponents agree with the two displayed
congruences.  Finally, the compatibility defect is a crossed
cocycle that vanishes on the arithmetic subgroup.  Matching
the multiplier of any full-GT element by an arithmetic element
reduces its value to the multiplier-one case proved here.
\end{proof}
